% Scheduler Self-Security — Article 1 (gate G8)
% Data: article1/paper/tables/scheduler_security_tests.csv (produced by
% article1/experiments/scheduler-security/run_scheduler_security_tests.sh via
% kubectl auth can-i against the scheduler ServiceAccount).

\section{Scheduler Self-Security}
\label{sec:scheduler-security}

Because the attestation-aware scheduler becomes a security-critical component ---
it decides placement and mints the placement token --- it must itself be
minimally privileged and fail-closed. We therefore treat the scheduler as part
of the attack surface and verify, on the live cluster, that its RBAC cannot be
abused beyond its role.

\subsection{Minimal RBAC}
The scheduler runs under a dedicated \texttt{ServiceAccount} whose ClusterRole
grants only: \emph{read} on \texttt{Pods}, \texttt{Nodes}, \texttt{Namespaces},
\texttt{ConfidentialInferencePolicy} and \texttt{AttestationEvidence};
\emph{create/update} on \texttt{AIPlacementDecision} (and its status);
\texttt{create} on \texttt{pods/binding} and \texttt{bindings} (needed by a
second scheduler); \texttt{create} on \texttt{Events}; and coordination
\texttt{Lease} access. It has \emph{no} wildcard verbs or resources.

\subsection{Fail-closed capability tests}
Table~\ref{tab:scheduler-security} reports the outcome of capability probes
(\texttt{kubectl auth can-i --as} the scheduler SA). The scheduler is confirmed
\emph{unable} to: create or update \texttt{AttestationEvidence} (S1--S2), create
\texttt{RawAttestationReport} (S3), modify \texttt{ConfidentialInferencePolicy}
(S4), patch \texttt{Node} labels (S5), read arbitrary \texttt{Secrets} (S6),
delete evidence (S10), create webhook configurations (S11), or use wildcard
permissions (S12). It \emph{is} able to read \texttt{AttestationEvidence} (S7),
create \texttt{AIPlacementDecision} (S8), and create \texttt{pods/binding} (S9)
--- exactly its role and nothing more.

\begin{table}[t]
  \centering
  \caption{Scheduler self-security probes (data:
  \texttt{tables/scheduler\_security\_tests.csv}). ``exp'' is the required
  outcome of \texttt{auth can-i}.}
  \label{tab:scheduler-security}
  \footnotesize
  \begin{tabular}{@{}llll@{}}
    \toprule
    \textbf{ID} & \textbf{Capability} & \textbf{exp} & \textbf{Property} \\
    \midrule
    S1  & create AttestationEvidence      & no  & no evidence forgery \\
    S2  & update AttestationEvidence      & no  & no evidence tamper \\
    S3  & create RawAttestationReport     & no  & not an attester \\
    S4  & update ConfidentialInferencePolicy & no & no policy tamper \\
    S5  & patch Nodes                     & no  & no node-label forgery \\
    S6  & get Secrets                     & no  & least secret access \\
    S7  & get AttestationEvidence         & yes & read-only appraisal \\
    S8  & create AIPlacementDecision      & yes & its output \\
    S9  & create pods/binding             & yes & second-scheduler bind \\
    S10 & delete AttestationEvidence      & no  & no evidence deletion \\
    S11 & create MutatingWebhookConfig    & no  & no admission hijack \\
    S12 & \texttt{*} \texttt{*}           & no  & no wildcard \\
    \bottomrule
  \end{tabular}
\end{table}

\subsection{Signing-key protection and fake-scheduler resistance}
The placement-token Ed25519 private key is held in a Kubernetes \texttt{Secret}
mounted only by the scheduler (or supplied via a sealed env var); the scheduler
RBAC does not grant blanket \texttt{Secrets} read (S6). Consequently a rogue or
``fake'' scheduler cannot mint tokens that \texttt{verify-placement} accepts:
without the signing key it can bind a pod (if it forged RBAC), but the resulting
\texttt{AIPlacementDecision} would carry no valid token and would fail
independent verification. This composes with the trust chain
(\S\ref{sec:threat}): a compromised node cannot forge evidence, and a rogue
scheduler cannot forge a verifiable placement.

% IP REVIEW REQUIRED BEFORE SUBMISSION
% (Signing-key handling and placement-token minting are described here.)
